\section{模型形状、推理效率与有效感受野}

达到某个 accuracy 后，下一步不是盲目放大所有维度。Transformer 的 depth、width、MLP width、head 数与 patch size 会改变优化、activation memory、并行效率和推理延迟。讲座用 shape、speed 与 receptive field 三组分析说明：同样参数或 compute 下，形状选择可以显著改变结果。

\subsection{Uniform scaling 并不总是最佳}

所谓 \term{model shape} 是参数如何分配到深度、hidden width、MLP width、heads 与 patch resolution。统一放大所有维度简单，却可能把预算投到边际收益较低的方向。图中不同颜色代表不同 scaling 方式，比较它们在 relative compute 下的平均性能。

\lecturefigure{slide-23-transformer-shape.jpg}{ViT 的 depth、width、patch size 与 MLP width 对 compute efficiency 的影响不同。}{官方视频 46:20。}

\begin{knowledgebox}{读图：曲线接近不代表系统等价}
纵轴是平均分，横轴是 relative compute。较高曲线表示同预算下更准确。Depth scaling 可能提高表达层级却增加串行路径；width 更利于并行但扩大 activation；更小 patch 提高空间分辨率却使 attention 成本激增。最终 shape 应同时考虑训练吞吐、推理 batch、memory 与目标分辨率。
\end{knowledgebox}

\subsection{Inference speed：复杂度常数和硬件映射会反转排名}

理论 FLOPs 只描述运算数量，真实 inference speed 还受 kernel implementation、batch size、memory bandwidth 与 accelerator shape 影响。ViT 的 dense matrix operations 容易利用现代硬件，但 token 数增加会迅速扩大 attention 与 activation。

\lecturefigure{slide-24-inference-speed.jpg}{不同 ResNet 与 ViT 变体的推理吞吐随图像规模/模型变化。}{官方视频 47:56。}

\begin{knowledgebox}{读图：速度曲线需要 hardware context}
纵轴代表推理速度或吞吐，横轴对应输入规模/模型配置。比较时必须固定 accelerator、precision、batch size、compiler 与 preprocessing。一个模型在大 batch 吞吐领先，可能在 batch-1 latency 落后；研究训练友好也不等于移动端友好。
\end{knowledgebox}

\begin{warningbox}{Patch size 是隐藏的系统旋钮}
减小 $P$ 能保留细粒度信息，却让 $N=HW/P^2$ 增长。它同时影响 attention FLOPs、KV/activation memory、数据加载后的 tensor shape 和下游 dense prediction resolution。报告“ViT-B”而不报告 `/16` 或 `/32`，会遗漏关键成本变量。
\end{warningbox}

\subsection{Effective receptive field：attention 也会形成层级结构}

ViT 每层理论上可全局交互，但实际 attention head 可能在浅层偏局部、深层扩大范围。\term{effective receptive field} 指模型对不同空间距离的实际依赖，而非理论可达范围。它能帮助判断局部结构是被硬编码还是从数据中涌现。

\lecturefigure{slide-25-receptive-field.jpg}{不同层与 attention heads 的平均注意距离展示 ViT 学到的局部到全局结构。}{官方视频 49:20。}

\begin{knowledgebox}{读图：点云表示多头异质性}
横轴是 layer depth，纵轴是平均 attention distance；每个点可对应一个 head。浅层存在局部 heads，深层多数 head 覆盖更远范围，说明模型形成层级。但 attention weight 不等同于 feature attribution，也不能单独证明某 head 的因果功能；需要 masking 或 intervention 实验验证。
\end{knowledgebox}

\teachervoice{老师在这一段不断回到“同样规模怎么定义”。参数量、FLOPs、实际吞吐与表示容量并不一致。实践中应把模型命名拆成 depth、width、patch、resolution 与 precision，避免用一个总参数数掩盖系统差异。}

\subsection{本章小结}

模型扩展是资源分配问题。Shape 决定计算如何进入 depth/width/token resolution，硬件决定 FLOPs 如何转化为速度，有效感受野则检查模型实际学到的空间层级。三者共同约束“更大 ViT”是否值得。

